\documentclass[11pt,a4paper]{article}
\usepackage[margin=2.5cm]{geometry}
\usepackage{amsmath,amssymb,amsthm}
\usepackage{booktabs}
\usepackage{array}
\usepackage{hyperref}
\usepackage{multirow}

\theoremstyle{plain}
\newtheorem{theorem}{Theorem}
\newtheorem{proposition}[theorem]{Proposition}
\newtheorem{corollary}[theorem]{Corollary}
\newtheorem{conjecture}[theorem]{Conjecture}

\theoremstyle{definition}
\newtheorem{definition}[theorem]{Definition}

\theoremstyle{remark}
\newtheorem{remark}[theorem]{Remark}
\newtheorem{observation}[theorem]{Observation}

\newcommand{\Z}{\mathbb{Z}}
\newcommand{\R}{\mathbb{R}}
\newcommand{\F}{\mathbb{F}}
\newcommand{\Aut}{\operatorname{Aut}}
\newcommand{\PSL}{\operatorname{PSL}}
\newcommand{\PGL}{\operatorname{PGL}}
\newcommand{\SU}{\operatorname{SU}}
\newcommand{\Sab}{\operatorname{Sab}}
\newcommand{\PSU}{\operatorname{PSU}}

\title{Spectral Synthesis:\\
  A Unified Survey of Eigenvalues Across the\\
  Cubic Arc-Transitive Graph Programme}
\author{Robin Langer}
\date{}

\begin{document}
\maketitle

\begin{abstract}
We present a unified spectral survey of 30+ graphs arising from the
programme of constructing cubic arc-transitive graphs as Sabidussi
coset graphs of groups of Lie type. The survey spans four families:
the Foster census (13 small cubic arc-transitive graphs), the
$\PSL(2,q)$ tower (7 graphs from $\SU(2)$ over finite fields),
the exceptional Lie group $G_2(2)$ (3 graphs of varying regularity),
and voltage covers via cubic residues (6 bipartite lifts).
Every graph in the census is Ramanujan for its degree,
a fact we trace to the Weil conjectures.
Cross-family comparison reveals universal structural patterns:
eigenvalues in small algebraic number fields,
Fiedler values decaying as $O(1/q)$, and eigenvalue multiplicities
governed by the representation theory of the acting group.
\end{abstract}


% ==========================================================================
\section{Introduction}
% ==========================================================================

The Sabidussi construction $\Sab(G, H, HaH)$ builds a graph from a group
$G$, a subgroup $H$, and a double coset $HaH$: the vertices are cosets
$G/H$ and two cosets $Hx$, $Hy$ are adjacent iff $xy^{-1} \in HaH$. When
$|HaH|/|H| = 3$, the result is a cubic (3-regular) graph; when $G$ acts
arc-transitively, the graph is arc-transitive.

This paper synthesises spectral data from a computational census of such
graphs, focusing on four families that exemplify the broader programme of
realising finite groups of Lie type as graph automorphisms.
Each family is treated in a companion paper; here we extract the
cross-cutting themes.

The four companion papers are:
\begin{enumerate}
\item \textbf{Foster census} --- the 13 smallest cubic arc-transitive
  graphs ($|V| \le 30$), all computable without group theory.
\item \textbf{$\PSL(2,q)$ family} --- cubic Sabidussi graphs from the
  $A_1$ Lie type, with stabilisers $H \in \{C_3, S_3, S_4\}$.
\item \textbf{$G_2(2)$ family} --- three graphs (cubic, 6-regular,
  7-regular) from the smallest exceptional Lie group.
\item \textbf{Voltage graphs} --- cyclic covers $\Z_p \to K_2$ with cubic
  residue voltages, connecting number theory to graph spectra.
\end{enumerate}


% ==========================================================================
\section{The universal Ramanujan phenomenon}
% ==========================================================================

\begin{definition}
A connected $d$-regular graph is \emph{Ramanujan} if every nontrivial
eigenvalue $\lambda$ (i.e., $|\lambda| < d$) satisfies
$|\lambda| \le 2\sqrt{d-1}$.
\end{definition}

\begin{theorem}[Computational]
\label{thm:all-ramanujan}
Every graph in the census is Ramanujan for its degree:
\begin{itemize}
\item All 13 Foster census graphs (cubic, bound $2\sqrt{2} \approx 2.828$).
\item All 7 $\PSL(2,q)$ Sabidussi graphs (cubic).
\item The Tutte 12-cage (cubic, $\lambda_2 = \sqrt{6} \approx 2.449$).
\item The $GH(2,2)$ collinearity graph (6-regular, $\lambda_2 = 3$
  vs.\ bound $2\sqrt{5} \approx 4.472$).
\item The Meinhold-72 graph (7-regular, $\lambda_2 = \sqrt{13} \approx 3.606$
  vs.\ bound $2\sqrt{6} \approx 4.899$).
\item All 6 cubic residue voltage lifts $\mathrm{CR}(p)$
  ($p = 7, 13, 19, 31, 37, 43$).
\end{itemize}
\end{theorem}

This universality is not a coincidence. For the $\PSL(2,q)$ family, the
Ramanujan property follows from Deligne's proof of the Weil conjectures:
the nontrivial eigenvalues of a Hecke operator on $\mathrm{GL}(2)$ over
a function field are bounded by $2\sqrt{q}$, and the graph eigenvalues
are evaluations of these operators at specific representations.
The Lubotzky--Phillips--Sarnak construction~\cite{LPS} made this
connection explicit for Cayley graphs of $\PSL(2,q)$.

For the voltage lifts, the eigenvalues are Gauss-type character sums
$f_j = \sum_{v \in V} \omega^{jv}$ where $V$ is the set of cubic residues
and $\omega = e^{2\pi i/p}$. The Weil bound gives
$|f_j| \le 2\sqrt{p}$ for nontrivial characters, which after normalisation
yields $\lambda_2 \le 2\sqrt{2}$.

\begin{observation}
The $G_2(2)$ graphs are Ramanujan despite coming from an exceptional Lie
group rather than a classical one. This suggests the Ramanujan property
may extend beyond the $A_1$ setting to all Chevalley groups, at least for
small field sizes.
\end{observation}


% ==========================================================================
\section{Master spectral table}
% ==========================================================================

Table~\ref{tab:master} collects the key spectral invariants across all
families. Graphs are grouped by family and sorted by vertex count within
each family.

\begin{table}[htbp]
\centering
\caption{Spectral invariants of all graphs in the census.
  $d$ = degree, $g$ = girth, $\lambda_2$ = second-largest adjacency eigenvalue,
  $\mu_1$ = Fiedler (algebraic connectivity), \#$\lambda$ = number of
  distinct eigenvalues. All are Ramanujan.}
\label{tab:master}
\smallskip
\renewcommand{\arraystretch}{1.05}
\small
\begin{tabular}{@{}l r r r r r r l@{}}
\toprule
Graph & $|V|$ & $d$ & $g$ & $\lambda_2$ & $\mu_1$ & \#$\lambda$ & Source \\
\midrule
\multicolumn{8}{@{}l}{\textit{Foster census}} \\
$K_4$         &   4 & 3 & 3 & $-1.000$ & $4.000$ &  2 & \\
$K_{3,3}$     &   6 & 3 & 4 & $0.000$  & $3.000$ &  3 & \\
Cube $Q_3$    &   8 & 3 & 4 & $1.000$  & $2.000$ &  4 & $\Z_2^3$ \\
Petersen      &  10 & 3 & 5 & $1.000$  & $2.000$ &  3 & $\PSL(2,5)/S_3$ \\
Heawood       &  14 & 3 & 6 & $\sqrt{2}$  & $1.586$ &  4 & $\PGL(2,7)/S_4$ \\
M\"ob.-Kantor &  16 & 3 & 6 & $\sqrt{3}$  & $1.268$ &  6 & \\
Pappus        &  18 & 3 & 6 & $\sqrt{3}$  & $1.268$ &  5 & \\
Dodecahedron  &  20 & 3 & 5 & $\sqrt{5}$  & $0.764$ &  6 & $A_5$ \\
Desargues     &  20 & 3 & 6 & $2.000$  & $1.000$ &  6 & \\
Nauru         &  24 & 3 & 6 & $2.000$  & $1.000$ &  7 & \\
F26A          &  26 & 3 & 6 & $2.074$  & $0.926$ &  6 & $\Z_{13}$ \\
Coxeter       &  28 & 3 & 7 & $2.000$  & $1.000$ &  5 & \\
Tutte-Coxeter &  30 & 3 & 8 & $2.000$  & $1.000$ &  5 & $\mathrm{GQ}(2,2)$ \\
\midrule
\multicolumn{8}{@{}l}{\textit{$\PSL(2,q)$ Sabidussi family ($A_1$ type)}} \\
$\PSL(2,5)/S_3$  &  10 & 3 & 5 & $1.000$  & $2.000$ &  3 & Petersen \\
$\PSL(2,5)/C_3$  &  20 & 3 & 5 & $\sqrt{5}$  & $0.764$ &  6 & $2 \times$ Petersen \\
$\PSL(2,7)/C_3$  &  56 & 3 & 6 & $\sqrt{2}{+}1$  & $0.586$ &  9 & Klein \\
$\PSL(2,17)/S_4$ & 102 & 3 & 9 & $2.562$  & $0.438$ &  8 & Biggs-Smith \\
$\PSL(2,13)/S_3$ & 182 & 3 & 7 & $2.675$  & $0.325$ & 16 & \\
$\PSL(2,11)/C_3$ & 220 & 3 &10 & $\varphi{+}1$ & $0.382$ & 19 & $\varphi = \tfrac{1+\sqrt{5}}{2}$ \\
$\PSL(2,17)/S_3$ & 408 & 3 & 9 & $2.709$  & $0.291$ & 24 & \\
\midrule
\multicolumn{8}{@{}l}{\textit{$G_2(2)$ exceptional family}} \\
GH(2,2) collin.  &  63 & 6 & 3 & $3.000$  & $3.000$ &  4 & $G_2(2)$ \\
Meinhold-72       &  72 & 7 & 4 & $\sqrt{13}$ & $3.394$ &  6 & $G_2(2)$ \\
Tutte 12-cage    & 126 & 3 &12 & $\sqrt{6}$  & $0.551$ &  7 & $G_2(2)$ \\
\midrule
\multicolumn{8}{@{}l}{\textit{Cubic residue voltage lifts $\mathrm{CR}(p)/K_2$}} \\
$\mathrm{CR}(7)$  &  14 & 3 & 6 & $\sqrt{2}$  & $1.586$ &  4 & $= $ Heawood \\
$\mathrm{CR}(13)$ &  26 & 3 & 6 & $2.074$  & $0.926$ &  6 & $= $ F26A \\
$\mathrm{CR}(19)$ &  38 & 3 & 6 & $2.347$  & $0.653$ &  8 & \\
$\mathrm{CR}(31)$ &  62 & 3 & 6 & $2.591$  & $0.409$ &  8 & \\
$\mathrm{CR}(37)$ &  74 & 3 & 6 & $2.655$  & $0.345$ &  8 & \\
$\mathrm{CR}(43)$ &  86 & 3 & 6 & $2.702$  & $0.298$ &  8 & \\
\bottomrule
\end{tabular}
\end{table}


% ==========================================================================
\section{Eigenvalues as algebraic numbers}
% ==========================================================================

A striking pattern across all families is that eigenvalues belong to small
algebraic number fields over $\mathbb{Q}$.

\begin{observation}[Algebraic number fields in the spectrum]
\label{obs:fields}
The following fields appear:
\begin{center}
\small
\begin{tabular}{@{}l l l@{}}
\toprule
Field & Graphs & Group-theoretic origin \\
\midrule
$\mathbb{Q}$ & $K_4$, $K_{3,3}$, Cube, Petersen, & Integer eigenvalues from \\
              & Coxeter, Tutte-Coxeter, GH(2,2) & highly symmetric quotients \\
$\mathbb{Q}(\sqrt{2})$ & Heawood, Klein, Biggs-Smith, & $|S_4| = 24 = 2^3 \cdot 3$; \\
              & Tutte 12-cage, CR(7) & $\sqrt{2}$ from $S_4$ representations \\
$\mathbb{Q}(\sqrt{3})$ & M\"obius-Kantor, Pappus & $\sqrt{3}$ from order-3 elements \\
$\mathbb{Q}(\sqrt{5})$ & Dodecahedron, PSL(2,5)/C$_3$, & $|A_5| = 60$; \\
              & PSL(2,11)/C$_3$ & golden ratio $\varphi \in \mathbb{Q}(\sqrt{5})$ \\
$\mathbb{Q}(\sqrt{6})$ & Tutte 12-cage & $12096 = 2^6 \cdot 3^3 \cdot 7$ \\
$\mathbb{Q}(\sqrt{13})$ & F26A, Meinhold-72, CR(13) & $\Z_{13}$ symmetry \\
\bottomrule
\end{tabular}
\end{center}
\end{observation}

The discriminant of the number field consistently matches the prime
structure of the acting group. When $p$ divides $|G|$ and appears as
the size of a cyclic factor in the stabiliser chain, $\sqrt{p}$ tends
to appear in the spectrum. This is not mysterious: the eigenvalues of a
Sabidussi graph are character sums over the double coset $HaH$, and
character values of finite groups lie in cyclotomic fields
$\mathbb{Q}(\zeta_n)$, whose real subfields contain square roots of the
relevant primes.

\begin{remark}
The golden ratio $\varphi = (1 + \sqrt{5})/2$ appears in three different
contexts: as $\lambda_2 = 1 = 2\varphi - 2$ of the Petersen graph
(trivially), as $\lambda_2 = \sqrt{5}$ of the dodecahedron ($A_5$ has
order divisible by 5), and as $\lambda_2 = \varphi + 1 \approx 2.618$
of $\PSL(2,11)/C_3$. Since $\PSL(2,11) \cong L_2(11)$ has order
$660 = 2^2 \cdot 3 \cdot 5 \cdot 11$, the appearance of $\sqrt{5}$ is
predicted by Observation~\ref{obs:fields}.
\end{remark}


% ==========================================================================
\section{Fiedler eigenvalue scaling}
% ==========================================================================

The Fiedler eigenvalue (algebraic connectivity) $\mu_1 = d - \lambda_2$
measures the bottleneck of the graph. For graphs from $\PSL(2,q)$, it
decays with~$q$.

\begin{proposition}[Fiedler scaling in the $A_1$ family]
For the $\PSL(2,q)/C_3$ subfamily, the Fiedler eigenvalue satisfies
$\mu_1 \approx C/q$ for a constant $C \approx 4$:
\begin{center}
\begin{tabular}{@{}c c c c@{}}
\toprule
$q$ & $|V|$ & $\mu_1$ & $q \cdot \mu_1$ \\
\midrule
5  &  20 & 0.764 & 3.82 \\
7  &  56 & 0.586 & 4.10 \\
11 & 220 & 0.382 & 4.20 \\
\bottomrule
\end{tabular}
\end{center}
The product $q \cdot \mu_1$ is approximately constant, suggesting
$\mu_1 \sim 4/q$ asymptotically.
\end{proposition}

This scaling has a representation-theoretic explanation. The Fiedler
eigenvalue corresponds to the eigenvalue of the adjacency operator on
the ``first nontrivial'' spherical function of the Hecke algebra
$\mathcal{H}(G, H)$. As $q$ grows, the Plancherel measure of
$\PSL(2,q)$ concentrates and this spherical eigenvalue approaches the
edge of the spectrum, driving $3 - \lambda_2 \to 0$.

For the voltage lift family $\mathrm{CR}(p)$, the scaling is
$\mu_1 \approx 6/p$:
\begin{center}
\begin{tabular}{@{}c c c c@{}}
\toprule
$p$ & $|V|$ & $\mu_1$ & $p \cdot \mu_1$ \\
\midrule
7  & 14 & 1.586 & 11.10 \\
13 & 26 & 0.926 & 12.03 \\
19 & 38 & 0.653 & 12.41 \\
31 & 62 & 0.409 & 12.69 \\
37 & 74 & 0.345 & 12.78 \\
43 & 86 & 0.298 & 12.82 \\
\bottomrule
\end{tabular}
\end{center}
Here $p \cdot \mu_1 \to 4\pi/\sqrt{3} \approx 12.87$ (conjectured),
arising from the equidistribution of cubic residues on the unit circle.


% ==========================================================================
\section{Spectral nesting and the tower of covers}
% ==========================================================================

The three $\PSL(2,q)$ families form a tower of covering maps:
\[
\frac{\PSL(2,q)}{C_3}
\;\xrightarrow{\ 2:1\ }\;
\frac{\PSL(2,q)}{S_3}
\;\xrightarrow{\ 4:1\ }\;
\frac{\PSL(2,q)}{S_4}
\]
corresponding to the subgroup chain $C_3 \subset S_3 \subset S_4$.
Each covering map induces spectral containment: the spectrum of the
quotient graph is a subset of the spectrum of the cover.

\begin{observation}[Spectral nesting at $q = 17$]
The Biggs-Smith graph $\PSL(2,17)/S_4$ has 8 distinct eigenvalues.
The larger graph $\PSL(2,17)/S_3$ has 24 distinct eigenvalues, and
every Biggs-Smith eigenvalue appears among them (with equal or greater
multiplicity). The 16 ``new'' eigenvalues in the $S_3$ graph arise from
representations of $\PSL(2,17)$ that are trivial on $S_4$ but nontrivial
on $S_3/S_4$.
\end{observation}

This nesting is a direct consequence of the Hecke algebra structure.
The spectrum of $\Sab(G, H, HaH)$ is determined by the spherical functions
of $\mathcal{H}(G, H)$. When $H' \supset H$, the Hecke algebra
$\mathcal{H}(G, H')$ embeds in $\mathcal{H}(G, H)$, and its spherical
functions restrict to spherical functions of the smaller algebra.


% ==========================================================================
\section{The Biggs-Smith extremum}
% ==========================================================================

Among all graphs in the census, the Biggs-Smith graph occupies a
distinguished position:

\begin{itemize}
\item It has the fewest distinct eigenvalues (8) relative to its size
  (102 vertices) --- fewer than the 56-vertex Klein graph (9 eigenvalues).
\item It is distance-transitive with diameter 7, achieving the minimum
  $\#\lambda = \text{diameter} + 1$.
\item Its stabiliser $S_4$ is \textbf{maximal} in $\PSL(2,17)$, leaving
  no room for intermediate quotients. This maximality forces
  distance-transitivity and collapses the spectrum.
\item Its $\lambda_2 \approx 2.562$ is lower than the larger
  $\PSL(2,13)/S_3$ ($\lambda_2 \approx 2.675$), despite having fewer
  vertices. The maximal stabiliser produces a ``tighter'' graph.
\end{itemize}

\begin{remark}
The Biggs-Smith graph is one of only 12 cubic distance-transitive graphs
(the full classification is due to Biggs and Smith, 1971). The other 11
all appear in the Foster census portion of our table. The spectral
simplicity of these graphs --- few eigenvalues, high multiplicities ---
is the spectral signature of distance-transitivity.
\end{remark}


% ==========================================================================
\section{$G_2(2)$: three faces of an exceptional group}
% ==========================================================================

The exceptional Lie group $G_2(2)$ (order 12096) provides a natural test
case for extending the programme beyond $A_1$. Three graphs share
$\Aut(G) = G_2(2)$ but have wildly different spectral profiles:

\begin{center}
\small
\begin{tabular}{@{}l c c l c@{}}
\toprule
Graph & $d$ & $\lambda_2$ & Spectrum shape & DT? \\
\midrule
Tutte 12-cage & 3 & $\sqrt{6}$ & $\{{\pm}3, {\pm}\sqrt{6}, {\pm}\sqrt{2}, 0\}$ & Yes \\
GH(2,2) collinearity & 6 & 3 & $\{6, 3, -1, -3\}$ & Yes \\
Meinhold-72 & 7 & $\sqrt{13}$ & $\{{\pm}7, {\pm}\sqrt{13}, {\pm}1\}$ & No \\
\bottomrule
\end{tabular}
\end{center}

\begin{observation}
The Tutte 12-cage and Meinhold-72 are both bipartite (spectrum symmetric
about~0), while the GH(2,2) collinearity graph is not. All three are
Ramanujan. The eigenvalue $\sqrt{6}$ in the Tutte 12-cage and $\sqrt{13}$
in Meinhold-72 both arise from the same $G_2(2)$ representation theory,
but evaluated on different double cosets.
\end{observation}

\begin{observation}
The GH(2,2) collinearity graph has spectral gap $6 - 3 = 3$, which equals
half the degree --- an unusually large gap. For comparison, the Petersen
graph (also strongly regular) has gap $3 - 1 = 2 = 2d/3$. This reflects the
exceptional expansion quality of $G_2$-geometries.
\end{observation}


% ==========================================================================
\section{Approaching the Alon--Boppana bound}
% ==========================================================================

The Alon--Boppana theorem states that for any infinite family of
$d$-regular graphs, $\liminf \lambda_2 \ge 2\sqrt{d-1}$.
The voltage lift family $\mathrm{CR}(p)$ provides a concrete example of
approach to this bound:

\[
\begin{array}{c|cccccc}
p & 7 & 13 & 19 & 31 & 37 & 43 \\
\hline
\lambda_2 & 1.414 & 2.074 & 2.347 & 2.591 & 2.655 & 2.702 \\
2\sqrt{2} - \lambda_2 & 1.414 & 0.754 & 0.482 & 0.238 & 0.174 & 0.127
\end{array}
\]

The approach is slow (roughly $O(1/\sqrt{p})$), consistent with the Weil
bound: $\lambda_2 \le 2\sqrt{2}\bigl(1 - c/\sqrt{p}\bigr)$ for some
constant $c$. These graphs are ``asymptotically optimal'' expanders in the
Ramanujan sense, saturating the bound from below.

The $\PSL(2,q)$ family also approaches the bound but more slowly,
because the graphs grow as $q^3$ (not linearly in $q$):
\[
\begin{array}{c|ccccc}
q & 5 & 7 & 11 & 13 & 17 \\
\hline
\lambda_2\text{ (best)} & 1.000 & 2.414 & 2.618 & 2.675 & 2.709 \\
2\sqrt{2} - \lambda_2 & 1.828 & 0.414 & 0.210 & 0.153 & 0.120
\end{array}
\]


% ==========================================================================
\section{The spectral--geometric dictionary}
% ==========================================================================

The census reveals a dictionary between spectral and geometric/algebraic
properties:

\begin{center}
\small
\begin{tabular}{@{}l l@{}}
\toprule
Spectral property & Geometric/algebraic meaning \\
\midrule
$\lambda_{\min} = -d$ & Bipartite (incidence graph) \\
$\#\lambda = \text{diam} + 1$ & Distance-transitive \\
$\lambda_2$ small & Large stabiliser, ``tight'' graph \\
$\mu_1$ large & Good bottleneck, small diameter \\
Integer eigenvalues & Highly degenerate double cosets \\
$\sqrt{p}$ in spectrum & $p \mid |\Aut(G)|$, character field $\mathbb{Q}(\sqrt{p})$ \\
Spectrum $\subset$ larger spectrum & Covering/quotient relationship \\
All Ramanujan & Group of Lie type (Weil conjectures) \\
\bottomrule
\end{tabular}
\end{center}

\begin{conjecture}[Universal Ramanujan property]
Every cubic arc-transitive graph arising as a Sabidussi construction
$\Sab(G, H, HaH)$ from a simple group $G$ of Lie type is Ramanujan.
More generally, every Sabidussi graph from a Chevalley group over
$\F_q$ is Ramanujan for its degree.
\end{conjecture}

The evidence is strong: all 30+ graphs tested satisfy the bound, spanning
the Lie types $A_1$ (PSL), $G_2$ (exceptional), and the abelian case
($\Z_p$ via Weil sums). The conjecture would follow from a suitable
generalisation of the Ramanujan--Petersson conjecture to automorphic forms
on arbitrary reductive groups over function fields --- a theorem of
Lafforgue for $\mathrm{GL}(n)$ but still open in general.


% ==========================================================================
\section{Open directions}
% ==========================================================================

\begin{enumerate}
\item \textbf{Other Lie types.} Compute spectra for Sabidussi graphs from
  $\PSL(3,q)$, $\mathrm{Sp}(4,q)$, $\PSU(3,q)$, and the Suzuki/Ree
  groups. The main obstacle is constructing cubic double cosets
  computationally.

\item \textbf{Fiedler asymptotics.} Prove the scaling
  $\mu_1 \sim C/q$ for the $\PSL(2,q)/C_3$ family and identify the
  constant $C$ in terms of the Plancherel measure.

\item \textbf{Eigenvalue field tower.} As $q$ grows, do the eigenvalue
  fields $\mathbb{Q}(\lambda_2)$ stabilise, grow, or cycle? For
  $\mathrm{CR}(p)$, the field is $\mathbb{Q}(\cos(2\pi/p))$, which
  grows with $p$. For $\PSL(2,q)$, the pattern is less clear.

\item \textbf{Spectral rigidity.} When does the spectrum of a Sabidussi
  graph determine the full automorphism group? The companion work on the
  Klein quartic (SPECTRAL\_RIGIDITY.md) shows that irreducibility of
  eigenspaces under the automorphism group prevents ``hidden symmetries.''
  Extend this to the full census.

\item \textbf{Halved cubes and quotients.} The hypercube $Q_n$ is Ramanujan
  for $n \le 6$. Can passage to the halved cube $\tfrac{1}{2}Q_n$ restore
  the Ramanujan property for $n \ge 7$?
\end{enumerate}


% ==========================================================================
\section{Conclusion}
% ==========================================================================

The universal Ramanujan property of Sabidussi graphs from groups of Lie
type is the central finding of this census. It is not a single theorem but
a web of connections: the Weil conjectures bound character sums, character
sums are graph eigenvalues, and the Ramanujan bound is the spectral shadow
of deep arithmetic geometry. The spectral data reveals this structure
concretely: eigenvalues in number fields predicted by group orders, Fiedler
values scaling with field size, and multiplicities reflecting representation
dimensions.

The programme of drawing Lie groups as graphs --- starting from SU(2) over
$\F_{17}$ producing the Biggs-Smith graph --- has a spectral counterpart:
\emph{hearing} the group in the graph's spectrum. The census shows that you
can hear quite a lot.

\begin{thebibliography}{10}
\bibitem{LPS}
A.~Lubotzky, R.~Phillips, P.~Sarnak,
\emph{Ramanujan graphs},
Combinatorica \textbf{8} (1988), 261--277.

\bibitem{Deligne}
P.~Deligne,
\emph{La conjecture de Weil. I},
Publ.\ Math.\ IH\'ES \textbf{43} (1974), 273--307.

\bibitem{Sabidussi}
G.~Sabidussi,
\emph{Vertex-transitive graphs},
Monatsh.\ Math.\ \textbf{68} (1964), 426--438.

\bibitem{BiggsSm}
N.~Biggs, D.~H.~Smith,
\emph{On trivalent graphs},
Bull.\ London Math.\ Soc.\ \textbf{3} (1971), 155--158.

\bibitem{BCN}
A.~E.~Brouwer, A.~M.~Cohen, A.~Neumaier,
\emph{Distance-Regular Graphs},
Springer, 1989.

\bibitem{Lafforgue}
L.~Lafforgue,
\emph{Chtoucas de Drinfeld et correspondance de Langlands},
Invent.\ Math.\ \textbf{147} (2002), 1--241.

\bibitem{Alon}
N.~Alon,
\emph{Eigenvalues and expanders},
Combinatorica \textbf{6} (1986), 83--96.

\bibitem{Weil}
A.~Weil,
\emph{Sur les courbes alg\'ebriques et les vari\'et\'es qui s'en d\'eduisent},
Hermann, 1948.
\end{thebibliography}

\end{document}
